\icode{E-EME-SI-L2-261121-004}
如图，\(M\)、\(N\)和\(P\)是以\(MN\)为直径的半圆弧上的三点，\(O\)点为半圆弧的圆心，\(\angle MOP =\) \,\qty{60}{\degree}。电荷量相等、符号相反的两个电荷分别置于\(M\)、\(N\)两点，这时\(P\)点电场强度的大小为\(E_1\)；若将\(N\)点处的点电荷移至\(P\)点，则\(P\)点的场强大小变为\(E_2\)。\(E_1\)与\(E_2\)之比为\paren[]

\begin{tikzpicture}[>=stealth,font=\scriptsize]
	\draw[dashed,cyan] (-1.5,0)node[above left,black]{\(M\)} arc(-180:0:1.5)node[above right,black]{\(N\)}--cycle (0,0)node[above,black]{\(O\)}--(-120:1.5)node[below left,black]{\(P\)};
	\draw[gray] (-.2,0) arc(-180:-120:.2);
	\node at(-.4,-.2){\qty{60}{\degree}};
	\foreach \c in {(-1.5,0),(0,0),(1.5,0),(-120:1.5)}
		\fill \c circle(.04);
\end{tikzpicture}
\begin{choices}
	\item \(1:2\)
	\item \(2:1\)
	\item \(2:\sqrt{3}\)
	\item \(4:\sqrt{3}\)
\end{choices}
\begin{answer}
	B\qquad 提示：
\end{answer}